\label{chap:83-11}

\section{Exponential identity of TRIT and coordination of the elliptic, parabolic, and hyperbolic regimes}

\subsection{Construction of the elliptic operator by means of the Paley matrix and the relation \texorpdfstring{\(J_{+1}^{2}=-I\)}{J² = -I}}

The root of \(-1\) is not introduced as a formal symbol without an
antecedent. The six units modulo nine are coordinatized over
\(\mathbb P^1(\mathbb F_5)\), and the quadratic character of \(\mathbb F_5\) determines the
Paley conference matrix. Its ternary reduction is
\[
 A_W=
 \begin{pmatrix}
 0&1&1&1&1&1\\
 1&0&1&2&2&1\\
 1&1&0&1&2&2\\
 1&2&1&0&1&2\\
 1&2&2&1&0&1\\
 1&1&2&2&1&0
 \end{pmatrix}\in M_6(\mathbb F_3).
\]
The conference identity reduces to
\begin{equation}\label{p12-83-c11-s1:eq:constantes-raiz-interna-menos-uno}
 \boxed{A_W^2=2I_6=-I_6.}
\end{equation}
Therefore, \(A_W\) is a quarter-turn operator in the ternary space.
After extending scalars to \(\mathbb F_9\), it decomposes into the conjugate
eigenspaces of the roots of \(X^2+1\). The common integral source
\[
 \mathcal Q_{\mathbb Z}=\mathbb Z[u]/(u^2+1)
\]
has a finite realization \(u\mapsto A_W\) and a real realization
\(u\mapsto J_{\rm e}\). The source coordinates the quadratic relation; it does not
identify fields of different characteristics.

\subsection{Unified family \texorpdfstring{\(\exp(tJ_\tau)=C_\tau(t)I+S_\tau(t)J_\tau\)}{exp(t J tau) = C tau(t) I + S tau(t) J tau}}

The three TRIT regimes are carried out by means of operators
\[
 J_{+1}^2=-I,\qquad J_0^2=0,\qquad J_{-1}^2=I.
\]
Separating the even and odd powers of the exponential yields a single family:
\begin{equation}\label{p12-83-c11-s1:eq:constantes-euler-extendido}
 \boxed{
 \exp(tJ_\tau)=C_\tau(t)I+S_\tau(t)J_\tau,
 \qquad \tau\in\{+1,0,-1\},}
\end{equation}
where
\[
 (C_{+1},S_{+1})=(\cos,\sin),\qquad
 (C_0,S_0)=(1,t),\qquad
 (C_{-1},S_{-1})=(\cosh,\sinh).
\]
\subsection{Elliptic, parabolic, and hyperbolic specializations and compatibility with the temporal reader}

Its three governing specializations are
\begin{equation}\label{p12-83-c11-s1:eq:constantes-euler-tres-ramas}
 \exp(\pi J_{+1})+I=0,\qquad
 \exp(J_0)=I+J_0,\qquad
 \exp\!\bigl(2(\log\varphi)J_{-1}\bigr)=F_{\rm av}^{\,2}.
\end{equation}
The elliptic branch contains the complex circular closure; the parabolic branch,
threshold transport; the hyperbolic branch, golden self-scaling. Euler’s equation
thus appears as a specialization of a family that coordinates rotation,
boundary and hyperbolic expansion.


\section{Hopf fibration, Villarceau circles, and Möbius holonomy as a topological realization of the elliptic regime}

\subsection{Complex structure and Hopf fibration}

Let \(S^3\subset\mathbb C^2\) and
\[
h(z_0,z_1)=
\bigl(2\Re(z_0\overline z_1),
2\Im(z_0\overline z_1),|z_0|^2-|z_1|^2\bigr)
\in S^2.
\]
On the Hopf torus
\[
T_\eta=\{(z_0,z_1):|z_0|=\cos\eta,
|z_1|=\sin\eta\},
\]
the curves
\[
\gamma_+(t)=(\cos\eta\,e^{it},\sin\eta\,e^{it}),
\]
\[
\gamma_-(t)=(\cos\eta\,e^{it},\sin\eta\,e^{-it})
\]
have distinct behaviors: \(\gamma_+\) is a Hopf fiber and
\(\gamma_-\) projects with degree two. Under stereographic projection, the
homology classes \((1,1)\) and \((1,-1)\) produce the two diagonal
Villarceau families.

The total complex conjugation
\[
K(z_0,z_1)=(\overline z_0,\overline z_1)
\]
preserves each unoriented family and reverses its orientation. It does not interchange
the two families. Partial conjugation can interchange them, but it is neither a
complex-linear map nor the standard global antiunitary. This correction
is essential to avoid manufacturing a particle--antiparticle symmetry from the
toroidal drawing.

\subsection{Double covering and spinor return}

The map \(SU(2)\to SO(3)\) is a double cover. A loop of
\(2\pi\) in \(SO(3)\) can lift to \(-I\) in \(SU(2)\); the loop of
\(4\pi\) recovers \(I\). The sign structure does not turn a circle into a
Möbius band. Obtaining a band requires a real line bundle
whose holonomy character is \(-1\).

Let \(L_\chi\to S^1\) be the associated bundle to
\[
\chi:\pi_1(S^1)\longrightarrow\{\pm1\},
\qquad \chi(1)=-1.
\]
Its total space is a Möbius band. The route \(C_M=-\operatorname{rev}\)
provides the sign candidate; the Villarceau circle provides the loop; the
physical intertwiner must prove that the former is the holonomy of the latter.

\subsection{Corrispondence between the spectrum \texorpdfstring{\(\{7,3\}\)}{{7,3}}, the ratio \texorpdfstring{\(5{:}2\)}{5:2} and the spinorial return}

With \(B_c=7I+2S\), define the oriented memory operator
\[
B_{\rm mem}^{\varepsilon}=5I+2\varepsilon S,
\qquad\varepsilon=\pm1.
\]
For \(\varepsilon=+1\),
\[
\Spec(B_{\rm mem}^{+})=\{7,3\},
\qquad\det B_{\rm mem}^{+}=21.
\]
Interpreting the values \((7,3)\) as the outer and inner equatorial boundaries
in a common positive chart yields
\[
R+r=7\ell,\qquad R-r=3\ell,
\]
and therefore
\[
R=5\ell,\qquad r=2\ell,\qquad
\sin\beta=\frac rR=\frac25.
\]
The Villarceau angle is
\[
\beta=\arcsin\frac25
=23.5781784782\ldots^\circ.
\]

The algebraic theorem \(\{7,3\}\leftrightarrow(5,2)\) is exact. The geometric
premise interpreting the spectrum as the two boundaries of a torus is a
declared interface. It must not be concealed, nor may the angle be used as a target.

The realization of the full block \(B_c\) retains another reading:
\(R:r=7:2\) and \(\beta=\arcsin(2/7)\). These are not rival values; they belong to
two functors: total block and memory mode. The physical priority of the latter
must be decided by the intertwiner with the dynamics, not by decimal proximity to
another constant.

\subsection{Angular momentum representation in the elliptic regime and phase space unit closure}

HMT’s native angle is a phase \(u\in\mathbb R/\mathbb Z\). The readers
publish
\[
\theta_{\rm rad}=2\pi u,
\qquad
\theta_{\rm deg}=360u.
\]
Radians and degrees are distinct charts of the same circular element. The radian is not more “ontological” by virtue of being dimensionless; its naturality follows from arc length. Degrees provide another absolute coordinate of the full turn. Any relation with \(\alpha^{-1}\), \(\varphi\), or CKM must first be formulated as an invariant quotient or holonomy and only thereafter evaluated in an angular chart.

The exact identity already available
\[
\tanh^2(2\log\varphi)=\frac59
\]
connects the Perron--Fibonacci channel to the spectral quotient of \(B_c\). Its
complement is \(4/9\), the normalized gap. This relation is structural.
It does not imply that every angle constructed with \(\varphi\) is a physical angle.

\subsection{Brouwer, Möbius, and degree}

Brouwer's fixed-point theorem, fractional linear Mobius transformations,
and the Mobius band are three distinct objects. They are related only
after fixing domain, action, and bundle. The elliptic advance and
retardation orbits may be modeled as two sections of a bundle with sign
holonomy; the Mobius transformation describes the conformal action in the chart;
Brouwer's theorem guarantees a fixed point for continuous maps of the disk.
None of these statements replaces the realization diagram.

\begin{fronteraBox}[title={Result and boundary}]
The degree-two character of the Hopf curve, the orientation of the
total conjugation, and the algebraic construction \(7,3\mapsto5,2\) are proved. The
toroidal reading (5{:}2), Möbius holonomy, and the
particle--antiparticle interpretation form a chain of interfaces that must share a
single intertwiner. The volume does not force that intertwiner through an
angular coincidence.
\end{fronteraBox}
